\centerline{\bf \Huge Тренировка региона ВсОШ}
\centerline{\bf \Huge 9 класс}

\problem{9.1}
Даны четыре последовательных натуральных числа, бо́льших 100. Докажите, что из них можно выбрать три числа, сумма которых представляется в виде произведения трёх различных натуральных чисел, бо́льших 1.

\problem{9.2}
На прямоугольном столе лежат несколько картонных прямоугольников. Их стороны параллельны сторонам стола. Размеры прямоугольников могут различаться, они могут перекрываться, но никакие два прямоугольника не могут иметь 4 общих вершины. Может ли оказаться, что каждая точка, являющаяся вершиной прямоугольника, является вершиной ровно трёх прямоугольников?

\problem{9.3}
Дан треугольник $A B C$. На внешней биссектрисе угла $A B C$ отмечена точка $D$, лежащая внутри угла $B A C$, такая, что $\angle B C D=60^{\circ}$. Известно, что $C D=2 A B$. Точка $M$ - середина отрезка $B D$. Докажите, что треугольник $A M C$ - равнобедренный.

\problem{9.4}
На доске нарисован выпуклый $n$-угольник ( $n \geqslant 4$ ). Каждую его вершину надо окрасить либо в чёрный, либо в белый цвет. Назовём диагональ разноцветной, если её концы окрашены в разные цвета. Раскраску вершин назовём хорошей, если $n$-угольник можно разбить на треугольники разноцветными диагоналями, не имеющими общих точек (кроме вершин). Найдите количество хороших раскрасок.

\problem{9.5}
Петя и Вася играют в следующую игру. Петя выбирает 100 (не обязательно различных) неотрицательных чисел $x_1, x_2$, $\ldots, x_{100}$, сумма которых равна 1 . Вася разбивает их на 50 пар по своему усмотрению, считает произведение чисел в каждой паре и выписывает на доску наибольшее из 50 полученных произведений. Петя хочет, чтобы число на доске оказалось как можно больше, а Вася - чтобы оно было как можно меньше. Какое число окажется на доске при правильной игре?